% GENERATED FILE. Edit canonical lessons, then run npm run generate:books.
% canonical-source-hash: fnv1a64:7109f4167680def2
% canonical-lessons: FR-C292-enqueteur, FR-C292-inspecteur, FR-C292-detective, FR-C292-indice, FR-C292-empreinte

\chapter{Les enquêteurs}
\label{ch:fr-les-enqueteurs}
\hlchaptermodality{292}

\begin{chapteropening}
\textbf{By the end of this chapter:} \emph{I can talk about the people who look into a crime and what they find: the investigator, the inspector, the detective, a clue and a print.}

Complete the last lesson of chapter 292: I can talk about the people who look into a crime and what they find: the investigator, the inspector, the detective, a clue and a print.
\end{chapteropening}

\section[l'enquêteur]{l'enquêteur — the investigator}

\label{lesson:FR-C292-enqueteur}



\begin{quote}
Before the new one: say the French for the patrol, then the French for the squad.
\end{quote}

\subsection*{You'll want to know: l'enquêteur}
\textbf{l'enquêteur} — \textquotedblleft{}the investigator\textquotedblright{}.

\begin{quote}
Les enquêteurs ont parlé avec les voisins.
\end{quote}

\begin{quote}
\emph{The investigators talked to the neighbours.}
\end{quote}

Feminine: \emph{l'enquêtrice}.

\begin{cousinweb}
\emph{enquêteur} is built on the noun for \textquotedblleft{}an inquiry\textquotedblright{}, which is inherited from Latin \emph{inquīsīta} \textquotedblleft{}things sought\textquotedblright{}, from \emph{inquīrere} \textquotedblleft{}to seek into\textquotedblright{}, built on \emph{quaerere} \textquotedblleft{}to seek\textquotedblright{} (Lewis \& Short).
\end{cousinweb}

\begin{grammarlens}[title={What you've built}]
One more word for telling, explaining and describing at B1.
\end{grammarlens}

\subsection*{Your turn}
\begin{itemize}
\raggedright
  \item \emph{Say it:} \emph{l'enquêteur}
  \item \emph{Say it:} \emph{l'enquêteur}, once more
  \item \emph{Say it:} a sentence of your own with \emph{l'enquêteur}
  \item \emph{Recall:} say the French for the patrol, then the French for the squad, then say \emph{l'enquêteur} again
\end{itemize}

\begin{tcolorbox}[breakable,colback=teal!4,colframe=teal!35!black,title={Before you move on}]
What does l'enquêteur mean? (\textquotedblleft{}The investigator\textquotedblright{}.) Say it once more.
\end{tcolorbox}
\section[l'inspecteur]{l'inspecteur — the inspector, the police inspector}

\label{lesson:FR-C292-inspecteur}



\begin{quote}
Before the new one: say the French for the squad, then the French for the investigator.
\end{quote}

\subsection*{You'll want to know: l'inspecteur}
\textbf{l'inspecteur} — \textquotedblleft{}the inspector, the police inspector\textquotedblright{}.

\begin{quote}
L'inspecteur a posé beaucoup de questions.
\end{quote}

\begin{quote}
\emph{The inspector asked a lot of questions.}
\end{quote}

Also at school or at work: \emph{un inspecteur des impôts} — a tax inspector.

\begin{cousinweb}
\emph{inspecteur} is a learned borrowing from Latin \emph{īnspector} \textquotedblleft{}examiner\textquotedblright{}, from \emph{īnspicere} \textquotedblleft{}to look into\textquotedblright{} (Lewis \& Short).
\end{cousinweb}

\begin{grammarlens}[title={What you've built}]
One more word for telling, explaining and describing at B1.
\end{grammarlens}

\subsection*{Your turn}
\begin{itemize}
\raggedright
  \item \emph{Say it:} \emph{l'inspecteur}
  \item \emph{Say it:} \emph{l'inspecteur}, once more
  \item \emph{Say it:} a sentence of your own with \emph{l'inspecteur}
  \item \emph{Recall:} say the French for the squad, then the French for the investigator, then say \emph{l'inspecteur} again
\end{itemize}

\begin{tcolorbox}[breakable,colback=teal!4,colframe=teal!35!black,title={Before you move on}]
What does l'inspecteur mean? (\textquotedblleft{}The inspector, the police inspector\textquotedblright{}.) Say it once more.
\end{tcolorbox}
\section[le détective]{le détective — the private detective}

\label{lesson:FR-C292-detective}



\begin{quote}
Before the new one: say the French for the investigator, then the French for the inspector.
\end{quote}

\subsection*{You'll want to know: le détective}
\textbf{le détective} — \textquotedblleft{}the private detective\textquotedblright{}.

\begin{quote}
Il lit des romans de détectives.
\end{quote}

\begin{quote}
\emph{He reads detective novels.}
\end{quote}

\emph{Un détective privé} works for clients; police officers are \emph{inspecteurs} or \emph{enquêteurs}.

\begin{cousinweb}
\emph{détective} is usually traced to English \emph{detective} (EtymDB lists no parent), which is built on Latin \emph{dētēctus}, from \emph{dētegere} \textquotedblleft{}to uncover\textquotedblright{} (Lewis \& Short).
\end{cousinweb}

\begin{grammarlens}[title={What you've built}]
One more word for telling, explaining and describing at B1.
\end{grammarlens}

\subsection*{Your turn}
\begin{itemize}
\raggedright
  \item \emph{Say it:} \emph{le détective}
  \item \emph{Say it:} \emph{le détective}, once more
  \item \emph{Say it:} a sentence of your own with \emph{le détective}
  \item \emph{Recall:} say the French for the investigator, then the French for the inspector, then say \emph{le détective} again
\end{itemize}

\begin{tcolorbox}[breakable,colback=teal!4,colframe=teal!35!black,title={Before you move on}]
What does le détective mean? (\textquotedblleft{}The private detective\textquotedblright{}.) Say it once more.
\end{tcolorbox}
\section[l'indice]{l'indice — the clue; the index, the rating}

\label{lesson:FR-C292-indice}



\begin{quote}
Before the new one: say the French for the inspector, then the French for the private detective.
\end{quote}

\subsection*{You'll want to know: l'indice}
\textbf{l'indice} — \textquotedblleft{}the clue; the index, the rating\textquotedblright{}.

\begin{quote}
Le détective a trouvé un indice important.
\end{quote}

\begin{quote}
\emph{The detective found an important clue.}
\end{quote}

Also a number on a scale: \emph{l'indice des prix}.

\begin{cousinweb}
\emph{indice} is a learned borrowing from Latin \emph{indicium} \textquotedblleft{}a sign, a proof\textquotedblright{} (EtymDB), built on \emph{index} \textquotedblleft{}pointer, sign\textquotedblright{}.
\end{cousinweb}

\begin{grammarlens}[title={What you've built}]
One more word for telling, explaining and describing at B1.
\end{grammarlens}

\subsection*{Your turn}
\begin{itemize}
\raggedright
  \item \emph{Say it:} \emph{l'indice}
  \item \emph{Say it:} \emph{l'indice}, once more
  \item \emph{Say it:} a sentence of your own with \emph{l'indice}
  \item \emph{Recall:} say the French for the inspector, then the French for the private detective, then say \emph{l'indice} again
\end{itemize}

\begin{tcolorbox}[breakable,colback=teal!4,colframe=teal!35!black,title={Before you move on}]
What does l'indice mean? (\textquotedblleft{}The clue; the index, the rating\textquotedblright{}.) Say it once more.
\end{tcolorbox}
\section[l'empreinte]{l'empreinte — the print, the fingerprint}

\label{lesson:FR-C292-empreinte}



\begin{quote}
Before the new one: say the French for the private detective, then the French for the clue.
\end{quote}

\subsection*{You'll want to know: l'empreinte}
\textbf{l'empreinte} — \textquotedblleft{}the print, the fingerprint\textquotedblright{}.

\begin{quote}
Il y avait des empreintes sur le verre.
\end{quote}

\begin{quote}
\emph{There were prints on the glass.}
\end{quote}

\emph{Les empreintes digitales} — fingerprints; \emph{l'empreinte carbone} — the carbon footprint.

\begin{cousinweb}
\emph{empreinte} is the past participle of Old French \emph{empreindre} \textquotedblleft{}to press in\textquotedblright{}, which is inherited from Latin \emph{imprimere} \textquotedblleft{}to press into\textquotedblright{} (Lewis \& Short).
\end{cousinweb}

\begin{grammarlens}[title={What you've built}]
One more word for telling, explaining and describing at B1.
\end{grammarlens}

\subsection*{Your turn}
\begin{itemize}
\raggedright
  \item \emph{Say it:} \emph{l'empreinte}
  \item \emph{Say it:} \emph{l'empreinte}, once more
  \item \emph{Say it:} a sentence of your own with \emph{l'empreinte}
  \item \emph{Recall:} say the French for the private detective, then the French for the clue, then say \emph{l'empreinte} again
\end{itemize}

\begin{tcolorbox}[breakable,colback=teal!4,colframe=teal!35!black,title={Before you move on}]
What does l'empreinte mean? (\textquotedblleft{}The print, the fingerprint\textquotedblright{}.) Say it once more.
\end{tcolorbox}
